% ============================================================
% Template LaTeX ABNT NBR 14724/2024 — IFB
% Persona e Pesquisa de Mercado — Luminária de Mesa
% Atividade 3 — Metodologia e Prática de Projeto
% ============================================================

\documentclass[
  12pt,
  a4paper,
  oneside,
  brazil
]{article}

% ── Codificação e Idioma ──────────────────────────────────────────────────────
\usepackage[T1]{fontenc}
\usepackage[utf8]{inputenc}
\usepackage[brazil]{babel}

% ── Fonte Times New Roman ─────────────────────────────────────────────────────
\usepackage{mathptmx}

% ── Geometria da página — margens ABNT anverso ────────────────────────────────
\usepackage[
  a4paper,
  left=3cm,
  top=3cm,
  right=2cm,
  bottom=2cm,
  headheight=14pt,
  headsep=0.5cm
]{geometry}

% ── Espaçamento 1,5 entre linhas (ABNT) ──────────────────────────────────────
\usepackage{setspace}
\onehalfspacing

% ── Recuo do parágrafo: 1,25cm (ABNT) ────────────────────────────────────────
\usepackage{indentfirst}
\setlength{\parindent}{1.25cm}
\setlength{\parskip}{0pt}

% ── Formatação de seções — numeração progressiva NBR 6024 ────────────────────
\usepackage{titlesec}

\titleformat{\section}
  {\normalfont\normalsize\bfseries\MakeUppercase}
  {\thesection\quad}
  {0em}{}

\titleformat{\subsection}
  {\normalfont\normalsize\bfseries}
  {\thesubsection\quad}
  {0em}{}

\titleformat{\subsubsection}
  {\normalfont\normalsize\bfseries\itshape}
  {\thesubsubsection\quad}
  {0em}{}

\titlespacing*{\section}      {0pt}{18pt}{12pt}
\titlespacing*{\subsection}   {0pt}{18pt}{6pt}
\titlespacing*{\subsubsection}{0pt}{12pt}{6pt}

% Seções primárias iniciam em nova página
\let\oldsection\section
\renewcommand{\section}{\clearpage\oldsection}

% ── Numeração de páginas — canto superior direito, 10pt ──────────────────────
\usepackage{fancyhdr}
\pagestyle{fancy}
\fancyhf{}
\fancyhead[R]{\fontsize{10}{12}\selectfont\thepage}
\renewcommand{\headrulewidth}{0pt}
\renewcommand{\footrulewidth}{0pt}

\fancypagestyle{plain}{
  \fancyhf{}
  \fancyhead[R]{\fontsize{10}{12}\selectfont\thepage}
  \renewcommand{\headrulewidth}{0pt}
}

% ── Listas ───────────────────────────────────────────────────────────────────
\usepackage{enumitem}
\setlist{nosep, leftmargin=*}
\providecommand{\tightlist}{\setlength{\itemsep}{0pt}\setlength{\parskip}{0pt}}

% ── Tabelas e Quadros ─────────────────────────────────────────────────────────
\usepackage{longtable}
\usepackage{booktabs}
\usepackage{array}
\usepackage{multirow}
\renewcommand{\arraystretch}{1.3}

% ── Legendas ─────────────────────────────────────────────────────────────────
\usepackage{caption}
\DeclareCaptionLabelSeparator{emdash}{\,---\,}
\DeclareCaptionFormat{abnt}{\fontsize{10}{12}\selectfont#1#2#3}
\captionsetup{
  format=abnt,
  justification=centering,
  singlelinecheck=false,
  labelsep=emdash,
  position=above,
  skip=3pt
}

% ── Sumário e Lista de Quadros ────────────────────────────────────────────────
\usepackage{tocloft}
\renewcommand{\contentsname}{SUMÁRIO}

\newlistof{quadros}{loq}{LISTA DE QUADROS}
\setlength{\cftquadrosindent}{0pt}
\setlength{\cftquadrosnumwidth}{0pt}
\makeatletter
\renewcommand{\listofquadros}{%
  \section*{LISTA DE QUADROS}%
  \@starttoc{loq}%
}
\makeatother

\newcommand{\quadrotitle}[2]{%
  \vspace{1em}%
  \phantomsection%
  \addcontentsline{loq}{quadros}{Quadro #1 --- #2}%
  \noindent\textbf{Quadro #1 --- #2}\par\vspace{0.3em}%
}
\newcommand{\quadrofonte}[1]{\par\noindent\textit{Fonte: #1}\vspace{1em}\par}

% ── URL e Hyperlinks ─────────────────────────────────────────────────────────
\usepackage{url}
\PassOptionsToPackage{hyphens}{url}
\usepackage[hidelinks, colorlinks=false, breaklinks=true]{hyperref}

% ── Figuras ──────────────────────────────────────────────────────────────────
\usepackage{graphicx}
\usepackage{float}

\date{}

% ============================================================
\begin{document}
% ============================================================

% ══════════════════════════════════════════════════════════════════════════════
% CAPA — NBR 14724/2024
% ══════════════════════════════════════════════════════════════════════════════
\begin{titlepage}
  \thispagestyle{empty}
  \begin{center}

    \includegraphics[width=0.25\textwidth]{../../../logo_ifb_despro.png}\\[1.8em]

    {\normalsize\bfseries
      Instituto Federal de Educação, Ciência e Tecnologia de Brasília --- IFB\\[0.3em]
      \textit{Campus} Samambaia\\[0.3em]
      Curso Superior de Tecnologia em Design de Produto
    }

    \vfill

    {\large\bfseries
      PERSONA E PESQUISA DE MERCADO --- LUMINÁRIA DE MESA
    }\\[1.2em]

    {\normalsize
      Atividade 3: Persona e Pesquisa de Mercado
    }

    \vfill

    {\normalsize Brasília/DF\\2026}

  \end{center}
\end{titlepage}
\setcounter{page}{2}

% ══════════════════════════════════════════════════════════════════════════════
% FOLHA DE ROSTO — NBR 14724/2024
% ══════════════════════════════════════════════════════════════════════════════
\begin{titlepage}
  \thispagestyle{empty}
  \begin{center}

    {\normalsize\bfseries VICTOR COSTA · ANNA JI · LUÍS COSTA}

    \vfill

    {\large\bfseries PERSONA E PESQUISA DE MERCADO --- LUMINÁRIA DE MESA}

    \vfill

    \begin{flushright}
      \begin{minipage}{0.5\textwidth}
        \normalsize
        Atividade apresentada ao Curso Superior de Tecnologia em Design de Produto
        do Instituto Federal de Educação, Ciência e Tecnologia de Brasília --
        Campus Samambaia, como requisito parcial para a aprovação na disciplina
        de Metodologia e Prática de Projeto.

        \vspace{1em}
        Professora: Profª. Fernanda Freitas Costa de Torres
      \end{minipage}
    \end{flushright}

    \vfill

    {\normalsize Brasília/DF -- 2026}

  \end{center}
\end{titlepage}

% ══════════════════════════════════════════════════════════════════════════════
% RESUMO — NBR 6028/2003
% ══════════════════════════════════════════════════════════════════════════════
\clearpage
\thispagestyle{empty}
\begin{center}
  {\normalsize\bfseries RESUMO}
\end{center}

\vspace{1em}
\noindent
Este documento apresenta a persona e a pesquisa de mercado para o projeto de
luminária de mesa desenvolvido na disciplina de Metodologia e Prática de Projeto.
A persona é construída a partir do perfil do público-alvo definido no briefing:
mulher de 40 a 45 anos, profissional liberal, moradora do Noroeste em Brasília,
com alto poder aquisitivo e consumo orientado por critérios estéticos e culturais.
A pesquisa de mercado mapeia o segmento de luminárias de design premium no Brasil,
identificando referências, faixas de preço e lacunas de mercado com implicações
diretas para o projeto.

\vspace{1em}
\noindent\textbf{Palavras-chave:} persona; pesquisa de mercado; luminária de mesa;
design autoral; mercado premium.

% ══════════════════════════════════════════════════════════════════════════════
% LISTA DE QUADROS
% ══════════════════════════════════════════════════════════════════════════════
\clearpage
\thispagestyle{empty}
\listofquadros

% ══════════════════════════════════════════════════════════════════════════════
% SUMÁRIO — NBR 14724/2024
% ══════════════════════════════════════════════════════════════════════════════
\clearpage
\thispagestyle{empty}
\tableofcontents

% ══════════════════════════════════════════════════════════════════════════════
% CONTEÚDO
% ══════════════════════════════════════════════════════════════════════════════

\section{Persona}

\subsection{Perfil Geral}

\noindent\textbf{Nome:} Beatriz Andrade\\[0.3em]
\noindent\textbf{Idade:} 43 anos\\[0.3em]
\noindent\textbf{Profissão:} Arquiteta sócia de escritório de pequeno porte
especializado em interiores residenciais\\[0.3em]
\noindent\textbf{Renda mensal:} R\$ 18.000--22.000\\[0.3em]
\noindent\textbf{Estado civil:} Divorciada, dois filhos adultos que não moram
com ela\\[0.3em]
\noindent\textbf{Moradia:} Apartamento de 110 m\textsuperscript{2} no Noroeste,
Brasília, comprado há quatro anos e reformado por ela mesma

\subsection{Estilo de Vida}

Beatriz acorda às 6h30, lê o jornal no tablet enquanto toma café, e chega ao
escritório por volta das 8h30. Trabalha com clientes exigentes e tem o hábito de
pesquisar referências de design continuamente, não como obrigação profissional, mas
como prazer cultivado desde a graduação. Suas férias são planejadas em torno de
cidades com arquitetura relevante: já foi a São Paulo para visitar a Bienal, a
Buenos Aires para ver museus e a Lisboa para caminhar pelos bairros históricos.

O apartamento dela é o laboratório permanente de seus próprios gostos. Ela troca
objetos com frequência, não por insatisfação compulsiva, mas porque encontra peças
melhores à medida que seu repertório cresce. A televisão foi abolida da sala de
estar há três anos: no lugar, uma estante de livros e um quadro de artista
brasiliense. O som ambiente é sempre algum jazz ou MPB instrumental.

Aos finais de semana, frequenta feiras de design, ateliês de cerâmica e livrarias
especializadas. Tem uma lista mental de objetos que ``vai comprar um dia'' e às
vezes esses objetos chegam de presente antes que ela os compre.

\subsection{Relação com Design e Objetos}

Beatriz compra poucos objetos, mas compra bem. A lógica que ela usa é simples:
prefere ter menos e ter melhor. Um objeto barato que ela troca em dois anos custa
mais, no final, do que um objeto caro que dura vinte.

Ela reconhece marcas de design sem precisar que elas se anunciem. Tem duas
luminárias Lumini no escritório, uma cadeira Zanine Caldas reproduzida em marceneiro
de confiança e um vaso de Takashi Murakami que ganhou de um cliente e que considera
``bonito demais para onde está, mas sem lugar melhor ainda''. Conhece Artemide,
Anglepoise e Flos, não porque as tem, mas porque sabe o que elas representam.

O que ela não tolera: objetos que parecem mais do que são. Acabamento que simula
material, laminado que imita madeira, plástico pintado de metal. Ela identifica isso
em segundos e a ilusão a incomoda mais do que o material simples honesto.

\subsection{Relação com a Luminária}

A luminária que ela receberia de presente provavelmente ficaria no quarto, sobre
o criado-mudo do lado esquerdo da cama, o lado em que ela lê antes de dormir. O
criado-mudo já tem uma luminária, mas é de um conjunto comprado na reforma e ela
nunca gostou completamente.

O que ela esperaria de uma luminária nova: que a luz fosse boa (quente, sem ofuscar),
que o objeto tivesse personalidade quando apagado, e que os materiais fossem reais.
Se for madeira, que seja madeira de verdade. Se for vidro, que se veja que foi feito
à mão.

O que a frustraria: descobrir que o difusor é plástico translúcido, que o cabo é
aquele fio branco padrão de luminária de loja de R\$ 49,90, ou que o ``latão'' é na
verdade alumínio pintado que descasca em seis meses.

A frase que a define: \textit{``Eu não compro objetos. Eu escolho companhia para
um espaço.''}

\section{Pesquisa de Mercado}

\subsection{Panorama do Segmento no Brasil}

O mercado de luminárias de design no Brasil pode ser dividido em três camadas: o
segmento popular (até R\$ 300, dominado por grandes varejistas e importação asiática),
o segmento intermediário (R\$ 300--800, com fabricantes nacionais e importados sem
posicionamento autoral forte) e o segmento premium/autoral (acima de R\$ 800, com
design assinado, materiais nobres e produção limitada).

O segmento premium é o menos consolidado no Brasil, mas vem crescendo desde meados
dos anos 2010, impulsionado pela expansão do mercado de interiores de alto padrão,
pelo interesse crescente em design nacional e pela maior circulação de referências
internacionais via redes sociais e publicações especializadas.

Brasília, em particular, tem um mercado consumidor de design bem desenvolvido para
uma cidade que não é o principal polo de produção do setor. A presença de profissionais
liberais e servidores públicos de alta remuneração, combinada com a cultura
arquitetônica da cidade, cria uma base de compradores sofisticados que precisam de
oferta à altura.

\subsection{Referências de Produtos e Marcas}

\noindent\textbf{Lumini (Brasil):} Uma das marcas mais reconhecidas de iluminação de
design no Brasil. Trabalha com linha residencial e arquitetural, com produtos a partir
de R\$ 600 e luminárias de mesa entre R\$ 1.200 e R\$ 3.500. Ponto forte: qualidade
técnica e distribuição nacional. Ponto fraco: linguagem formal mais próxima do
minimalismo industrial do que do design autoral com materialidade artesanal.

\noindent\textbf{Dpot (Brasil):} Multimarca com curadoria de design nacional e
internacional. Distribui peças de autores brasileiros e representa marcas europeias.
Referência de ambiente para o público-alvo, mas com preços que colocam a luminária
de mesa entre R\$ 2.000 e R\$ 8.000, teto acima do posicionamento deste projeto.

\noindent\textbf{Artemide (Itália):} Marca de referência mundial. Luminárias de mesa
a partir de R\$ 2.500 no mercado brasileiro. Design icônico, qualidade técnica
incomparável. Funciona mais como referência estética do que como concorrente direto,
dado o diferencial de preço e a ausência de identidade brasileira.

\noindent\textbf{Anglepoise (Reino Unido):} Reconhecida pela lamp Type 75 e seus
derivados. Preço no Brasil entre R\$ 1.800 e R\$ 3.200. Forte apelo ao público com
gosto clássico-contemporâneo. Não tem identidade de material artesanal, é
essencialmente metal industrial, o que a diferencia do posicionamento deste projeto.

\noindent\textbf{Produtores independentes brasileiros:} Há um conjunto crescente de
ateliês e designers independentes que produzem luminárias com materiais naturais,
madeira, cerâmica, couro, rattan, com preços entre R\$ 700 e R\$ 2.000 e distribuição
principalmente via Instagram e feiras de design. Esse é o segmento mais próximo do
projeto em termos de linguagem e processo, e o menos consolidado em termos de marca
e alcance.

\subsection{Análise de Preço}

O Quadro 1 apresenta a segmentação do mercado de luminárias de design no Brasil por
faixa de preço e perfil de produto.

\quadrotitle{1}{Segmentação do mercado de luminárias de design no Brasil}

\begin{longtable}[]{@{}
  >{\raggedright\arraybackslash}p{\dimexpr(\linewidth - 4\tabcolsep)*20/100\relax}
  >{\raggedright\arraybackslash}p{\dimexpr(\linewidth - 4\tabcolsep)*25/100\relax}
  >{\raggedright\arraybackslash}p{\dimexpr(\linewidth - 4\tabcolsep)*55/100\relax}
  @{}}
\toprule\noalign{}
\begin{minipage}[b]{\linewidth}\raggedright
\textbf{Segmento}
\end{minipage} &
\begin{minipage}[b]{\linewidth}\raggedright
\textbf{Faixa de preço}
\end{minipage} &
\begin{minipage}[b]{\linewidth}\raggedright
\textbf{Perfil de produto}
\end{minipage} \\
\midrule\noalign{}
\endhead
\bottomrule\noalign{}
\endlastfoot
Popular & até R\$ 300 & Importado, plástico, sem identidade autoral \\
Intermediário & R\$ 300--800 & Nacional ou importado, materiais mistos, design
genérico \\
Premium nacional & R\$ 800--2.000 & Design autoral, materiais nobres, pequena série \\
Premium importado & R\$ 2.000--8.000+ & Marcas internacionais, design icônico \\
\end{longtable}

\quadrofonte{elaborado pelos autores.}

O posicionamento-alvo deste projeto é \textbf{premium nacional}, entre R\$ 900 e
R\$ 1.800, com diferenciação por identidade cultural brasileira e materialidade
artesanal.

\subsection{Lacunas Identificadas}

A pesquisa aponta três lacunas que o projeto pode explorar:

\noindent\textbf{1. Identidade cultural explícita sem ser folclórica.} Produtos
brasileiros no segmento autoral frequentemente oscilam entre o minimalismo genérico
(que poderia ser de qualquer país) e o regionalismo excessivo (com elementos que
remetem mais ao artesanato popular do que ao design contemporâneo). Há espaço para
uma linguagem que dialogue com a arquitetura modernista de Brasília, sofisticada,
geométrica e ao mesmo tempo com presença, sem cair no pastiche.

\vspace{0.5em}
\noindent\textbf{2. Materiais artesanais com acabamento de nível premium.} Muitos
produtos independentes com boa proposta de material têm execução que não sustenta o
preço: vedações irregulares, cabos mal fixados, difusores com variação de espessura
visível. O público-alvo paga mais, mas exige que o objeto entregue o que promete ao
toque, não só na fotografia.

\vspace{0.5em}
\noindent\textbf{3. Objeto funcional nos dois estados.} Grande parte das luminárias
autorais brasileiras são interessantes acesas, mas discretas quando apagadas. A dupla
função, objeto decorativo quando desligado e fonte de luz quando ligado, é pouco
explorada no segmento e altamente valorizada pelo público que decora com intenção.

\subsection{Implicações para o Projeto}

A pesquisa confirma o posicionamento definido no briefing e adiciona três implicações:

\begin{itemize}
\tightlist
\item O acabamento precisa ser executado com rigor técnico equivalente ao de marcas
  consolidadas, mesmo que a produção seja artesanal de pequeno porte. A diferença
  entre R\$ 900 e R\$ 400 está principalmente na execução, não no conceito.
\item A comunicação do produto (embalagem, etiquetas, eventual catálogo) é parte
  inseparável do valor percebido no contexto de presente. Um objeto bem embalado e
  bem documentado se vende como presente mesmo sem o comprador ter visto o produto
  ao vivo.
\item A identidade cultural deve ser legível para quem tem repertório, mas não precisa
  de explicação. O freijó precisa ser identificável como madeira brasileira para quem
  conhece, e belo para quem não conhece.
\end{itemize}

\section{Referências}\label{referencias}

BÜRDEK, Bernhard E. \textbf{Design}: história, teoria e prática do design de produtos.
São Paulo: Edgard Blücher, 2006.

DESIGN BRASIL. \textbf{Panorama do design brasileiro 2024}. São Paulo: Museu da Casa
Brasileira, 2024.

MORAES, Dijon de. \textbf{Metaprojeto}: o design do design. São Paulo: Edgard Blücher,
2010.

NORMAN, Donald A. \textbf{Design emocional}: por que adoramos (ou detestamos) os objetos
do dia a dia. Rio de Janeiro: Rocco, 2008.

\end{document}
